\section{误差分解与可靠性检验}
\label{sec:validation}
\subsection{反射率拟合误差与厚度误差的区别}
实测数据没有给出厚度真值，因此不能计算厚度的实测相对误差，也不能把反射率预测误差写成厚度测量精度。设保留的反射率为$y_i$、预测为$\widehat y_i$、残差$e_i=y_i-\widehat y_i$，我们报告
\begin{align}
 \mathrm{RMSE}&=\sqrt{\frac{1}{N}\sum_{i=1}^N e_i^2}, &
 \mathrm{MAE}&=\frac{1}{N}\sum_{i=1}^N |e_i|,\\
 E_{\max}&=\max_i |e_i|, &
 R^2&=1-\frac{\sum_i e_i^2}{\sum_j\sum_{i\in j}(y_i-\overline y_j)^2}.
 \label{eq:validation-metrics}
\end{align}
其中$j$区分两种角度，各自扣除均值。RMSE、MAE和最大绝对残差的单位均为百分点。$R^2$反映模型对该窗口反射率变化的解释程度，不能检验折射率先验是否正确。

\begin{table}[htbp]\centering
\caption{在各自主分析波段上的反射率拟合指标}\label{tab:metrics}
\begin{tabular}{lrrrr}
\toprule
材料与模型 & RMSE/百分点 & MAE/百分点 & $R^2$ & 最大绝对残差\\\midrule
SiC两光束 & 0.02440 & 0.01744 & 0.997640 & 0.2381\\
SiC有效多光束核 & 0.02433 & 0.01738 & 0.997654 & 0.2348\\
Si两光束 & 4.36318 & 2.39116 & 0.890717 & 19.8474\\
Si多光束 & 0.67894 & 0.43700 & 0.997354 & 2.9360\\
\bottomrule
\end{tabular}\end{table}

\tabref{tab:metrics}应在材料内部、相同窗口与同一数据集上比较。硅模型的误差大于碳化硅模型，不意味着硅厚度必然测得更差，因为前者拟合更宽波段、更大的反射率变化和更强的界面响应。硅两束与多束的比较则具有相同数据基础，多光束模型在保持相同自由参数数量的情况下显著降低了误差。

\subsection{连续波数块的预测验证}
若逐点随机划分训练与验证集，相邻波数处的高度相关结构可能同时落入两组，验证结果会过于乐观。本文把碳化硅光谱按$50\cmiv$、硅光谱按$100\cmiv$划分连续块，将块编号对5取余分配为五折；同一块在两个角度下同时留出。模型在训练块上重新估计背景、标定和非线性参数，再预测留出块。

\begin{table}[htbp]\centering
\caption{五折连续波数块留出验证}\label{tab:cv}
\begin{tabular}{lrrrr}
\toprule
折次 & SiC两束 & SiC多束核 & Si两束 & Si多束\\\midrule
1 & 0.04553 & 0.04680 & 6.52455 & 0.86411\\
2 & 0.03021 & 0.02980 & 5.79395 & 0.64575\\
3 & 0.02751 & 0.02899 & 4.79589 & 0.49959\\
4 & 0.02640 & 0.02882 & 5.90708 & 0.58381\\
5 & 0.02010 & 0.02019 & 3.36228 & 0.83838\\
\bottomrule
\end{tabular}\end{table}

\tabref{tab:cv}中所有误差均以百分点计。碳化硅两光束模型合并留出RMSE为0.03202，多光束核为0.03307，增加约3.25\%；因此，没有足够预测收益支持用高阶形状参数替换主模型。硅的多光束模型在五个折次中均更优，36个连续块中34块降低了均方误差，支持保留高阶反射项。

本验证在全数据初步确定的条纹分支上热启动，检验的是该分支内重新拟合后的局部预测稳定性，不是每折从头独立选择条纹级次的验证。波段、模型结构及屏蔽段也已经通过同一批数据确定，所以这里报告为内部稳健性证据，不冒充外部样本验证。

基于独立残差假设的BIC可写成$N\log(\mathrm{SSE}/N)+k\log N$\cite{Schwarz1978}，但密集光谱点不满足简单独立重复的解释。碳化硅有效多光束核的原样本BIC有所改善，而分块预测变差，说明只按名义样本数惩罚复杂度仍可能选入弱小的系统残差。本文因此优先报告分块预测及厚度变化，而不以单一BIC数值给出物理存在性的结论。

\subsection{折射率、角度与窗口敏感性}
定义无量纲局部灵敏度
\begin{equation}
 S_p\simeq\frac{\widehat d(p+\Delta p)-\widehat d(p-\Delta p)}{2\Delta p}
             \frac{p}{\widehat d(p)}.
 \label{eq:normalized-sensitivity}
\end{equation}
对角度接近零或无自然比例尺度的参数，直接报告绝对扰动及厚度变化，避免无量纲化造成歧义。对于折射率整体缩放$n\mapsto sn$，重新求解所有可调参数，而非只将最终厚度乘以一个近似比例。

\paperfigure{figures/index_sensitivity.png}{折射率缩放对条件厚度估计的影响}{0.78}{fig:index-sensitivity}
\begin{table}[htbp]\centering
\caption{折射率整体缩放的情景比较}\label{tab:index-sensitivity}
\begin{tabular}{lrrr}
\toprule
材料 & 折射率缩放 & 厚度/$\mu$m & 相对主值变化/\%\\\midrule
SiC & 0.990 & 7.45942 & +1.017\\
SiC & 1.010 & 7.31071 & -0.997\\
Si & 0.990 & 3.44039 & +1.033\\
Si & 0.995 & 3.42271 & +0.514\\
Si & 1.000 & 3.40521 & -0.000\\
Si & 1.005 & 3.38790 & -0.509\\
Si & 1.010 & 3.37075 & -1.012\\
\bottomrule
\end{tabular}\end{table}

\figref{fig:index-sensitivity}与\tabref{tab:index-sensitivity}显示，折射率增大时厚度减小，与相位积的补偿关系一致。碳化硅在$n$整体改变$\pm1\%$时，厚度变为约7.31071--7.45942$\um$；硅对应约3.37075--3.44039$\um$。这些范围是人为设定的参数情景，不是材料折射率的不确定度测量，更不是厚度的95\%置信区间。两种材料的尺度响应说明，补充独立折射率信息比增加密集采样点更有助于解释绝对厚度。

硅两角共同偏移$-0.2^\circ$、$0^\circ$和$+0.2^\circ$时，共同厚度为3.40499、3.40521和3.40544$\um$。在这个假设扰动下，角度因素的影响小于$\pm1\%$折射率变化；该比较并不证明真实仪器角度误差就是$0.2^\circ$。

\begin{table}[htbp]\centering
\caption{有效波段内的窗口敏感性}\label{tab:windows}
\begin{tabular}{lrrr}
\toprule
材料 & 窗口/$\mathrm{cm}^{-1}$ & 厚度/$\mu$m & RMSE/百分点\\\midrule
SiC & 2000--3200 & 7.38098 & 0.02416\\
SiC & 2100--3300 & 7.39770 & 0.02127\\
SiC & 2200--3300 & 7.39546 & 0.02089\\
SiC & 2000--3000 & 7.37384 & 0.02178\\
SiC & 2000--3100 & 7.37977 & 0.02302\\
Si & 500--3900 & 3.40579 & 0.68159\\
Si & 600--3900 & 3.40694 & 0.64589\\
Si & 800--3900 & 3.40822 & 0.63946\\
Si & 1000--3900 & 3.40960 & 0.49216\\
Si & 1500--3900 & 3.40248 & 0.29989\\
\bottomrule
\end{tabular}\end{table}

\tabref{tab:windows}只列位于所用文献关系适用域内的窗口。碳化硅厚度范围为7.37384--7.39770$\um$，硅为3.40248--3.40960$\um$。改变窗口时有效数据数量不同，不能据此跨行比较BIC的绝对大小。碳化硅若保留$2300$--$2400\cmiv$异常段，得到7.39890$\um$，RMSE增至0.06282个百分点；厚度变化约0.20\%，但局部误差明显增大，说明排除该段主要改善了相位估计周围的异常影响。

\subsection{三因素联合情景比较}
单因素比较固定了其余条件，不能说明多个假设同时变化时的表现。
为此，将折射率比例、两角共同偏移与窗口选择各取三个水平，
按$L_9(3^3)$正交表安排九组情景，每组均重新估计原主模型中的辅助参数。
折射率水平为0.99、1.00、1.01，共同角偏移为$-0.2^\circ$、$0^\circ$和$+0.2^\circ$，
窗口选择及数值结果见表\ref{tab:joint-sensitivity}。
\begin{table}[htbp]
\centering
\caption{三因素三水平 L9 联合情景敏感性；变化均相对于各材料主模型}
\label{tab:joint-sensitivity}
\small
\setlength{\tabcolsep}{4pt}
\begin{tabular}{crrrrrrr}
\toprule
情景 & $s_n$ & $\Delta\theta/{}^\circ$ & 窗口级别 & $d_{\rm SiC}/\mu\mathrm m$ & 变化/\% & $d_{\rm Si}/\mu\mathrm m$ & 变化/\%\\
\midrule
1 & 0.99 & -0.2 & 1 & 7.45404 & +0.944 & 3.44021 & +1.028 \\
2 & 0.99 & +0.0 & 2 & 7.45942 & +1.017 & 3.44039 & +1.033 \\
3 & 0.99 & +0.2 & 3 & 7.47374 & +1.211 & 3.44236 & +1.091 \\
4 & 1.00 & -0.2 & 2 & 7.38355 & -0.010 & 3.40499 & -0.006 \\
5 & 1.00 & +0.0 & 3 & 7.39770 & +0.181 & 3.40692 & +0.050 \\
6 & 1.00 & +0.2 & 1 & 7.38053 & -0.051 & 3.40547 & +0.008 \\
7 & 1.01 & -0.2 & 3 & 7.32323 & -0.827 & 3.37223 & -0.969 \\
8 & 1.01 & +0.0 & 1 & 7.30621 & -1.058 & 3.37079 & -1.011 \\
9 & 1.01 & +0.2 & 2 & 7.31146 & -0.987 & 3.37097 & -1.006 \\
\bottomrule
\end{tabular}
\par\smallskip
\begin{tabular}{lrrrr}
\toprule
因素 & SiC 极差/$\mu\mathrm m$ & SiC 极差/\% & Si 极差/$\mu\mathrm m$ & Si 极差/\%\\
\midrule
折射率比例 & 0.148773 & 2.0147 & 0.069657 & 2.0456 \\
共同角度偏移 & 0.001636 & 0.0222 & 0.000458 & 0.0134 \\
波段选择 & 0.017961 & 0.2432 & 0.001717 & 0.0504 \\
\bottomrule
\end{tabular}
\par\smallskip
\begin{minipage}{0.98\textwidth}\footnotesize
窗口级别1、2、3对SiC分别表示$2000$--$3100$、$2000$--$3300$、$2100$--$3300$；对Si分别表示$450$--$3500$、$450$--$4000$、$600$--$4000$，单位均为$\mathrm{cm}^{-1}$。SiC各情景均剔除$2300$--$2400$。$s_n$同时乘入同一参照折射率函数；两个角度均增加$\Delta\theta$。SiC固定两束模型，Si固定含Drude衬底及有效往返包络的多束模型，每个情景重新估计原模型中的其余参数。
因素极差为该因素三个水平的平均厚度之最大值减最小值；百分比以各材料主厚度归一化。这些人为指定的范围不代表概率分布或置信区间。L9中交互效应可能与主效应混杂，该表不能独立识别交互作用、因果贡献或全局Sobol指数。
\end{minipage}
\end{table}


联合情景下，碳化硅厚度覆盖约$7.30621$--$7.47374\um$，硅覆盖约$3.37079$--$3.44236\um$。
按三个水平的平均厚度计算，折射率因素的极差分别为0.148773和0.069657$\um$，
大于本组窗口与角度因素的相应极差，延续了单因素检验中折射率尺度占主导的观察。
该设计只覆盖全部27种组合中的9种，并未独立分离交互效应，
因此不能把这些极差解释成因果贡献率或全局灵敏度指数。
所选水平也不是实测概率分布，联合范围应作为模型条件变化的情景结果报告。

\subsection{分角一致性与条件重抽样}
\paperfigure{figures/angle_consistency.png}{共同厚度与两个角度分别拟合的比较}{0.78}{fig:angle-consistency}
碳化硅两个角度独立给出7.40781和7.36078$\um$，差0.04703$\um$；硅独立给出3.42055和3.39283$\um$，差0.02772$\um$。\figref{fig:angle-consistency}中的差异均大于条件数值拟合的细小波动。同片晶圆不保证不同光斑处厚度完全相同，角度改变也可能伴随标定和衬底响应差异；目前缺少信息区分这些原因。共同厚度应解释为在共享参数约束下对两条谱线的折中估计。

为估计固定模型下残差结构对厚度拟合的影响，采用移动块残差重抽样思想\cite{Kunsch1989}。设第$b$次从残差序列抽到索引$I_b$，构造
\begin{equation}
 y_i^{*(b)}=\widehat y_i+(e_{I_b(i)}-\overline e),\qquad
 \widehat d^{*(b)}=\arg\min_d\min_{\boldsymbol\eta}
     \sum_i\left[y_i^{*(b)}-R_i(d,\boldsymbol\eta)\right]^2.
 \label{eq:block-bootstrap}
\end{equation}
两角使用相同索引块，保留它们之间的残差依赖。碳化硅每5点抽取1点，名义块宽约$50\cmiv$；硅每3点抽取1点，块宽约$144.6\cmiv$。每种材料重复120次，随机种子与原始样本顺序固定。

\paperfigure{figures/bootstrap.png}{固定模型与折射率条件下的厚度重抽样分布}{0.78}{fig:bootstrap}
\begin{table}[htbp]\centering
\caption{模型固定条件下的移动块残差重抽样}\label{tab:bootstrap}
\begin{tabular}{lrrrr}
\toprule
材料 & 次数 & 块宽/$\mathrm{cm}^{-1}$ & 厚度标准差/$\mu$m & 95\%分位范围/$\mu$m\\\midrule
SiC & 120 & 约50 & 0.005609 & [7.37480, 7.39474]\\
Si & 120 & 144.6 & 0.001212 & [3.40306, 3.40731]\\
\bottomrule
\end{tabular}\end{table}

\figref{fig:bootstrap}和\tabref{tab:bootstrap}描述的是已选模型条件下的拟合波动。循环块在窗口首尾连接，碳化硅序列还存在屏蔽段，名义块宽只是近似；局部残差也不一定严格平稳。故这些区间仅作为条件重抽样范围，不覆盖晶型、折射率、测点差异和仪器模型误差，不能作为完整计量不确定度。分角差异与折射率敏感性必须和它们一起报告\cite{JCGM2008}。

\subsection{代码的物理极限检查}
计算实现另外接受三种非实测的数值自检：将闭式多光束振幅与80次显式往返叠加比较；令膜与衬底折射率相同，检验厚度不再产生干涉响应；用给定参数生成无噪声数值光谱，再检查程序能否恢复输入参数。这些算例只验证公式和程序的实现，没有被用来替代附件结果。表\ref{tab:physical-checks}给出核对记录。
\begin{table}[htbp]\centering
\caption{用于代码核对的物理极限与无噪声算例}\label{tab:physical-checks}
\begin{tabular}{lr}
\toprule
检查内容 & 误差量\\\midrule
闭式振幅与80束显式叠加 & $6.66\times10^{-16}$\\
零界面对比度时的厚度响应 & 0\\
无噪声参数恢复最大相对误差 & $1.80\times10^{-15}$\\
\bottomrule
\end{tabular}\end{table}


硅闭式振幅与显式叠加的最大差为$6.66\times10^{-16}$，零反差情形的厚度响应差为0，无噪声参数恢复的最大相对误差为$1.80\times10^{-15}$。这些结果说明代数和单位实现相容，不能说明真实样品满足全部建模假设。非线性求解采用SciPy实现\cite{Virtanen2020}；多起点只能降低局部极小值风险，不构成全局最优性证明。
